%-------------------------------------------------------------------------------
% Mattia Giovannini — CV Europass in italiano (versione pubblica)
% Classe europasscv (CTAN, LPPL 1.3). Compilare con:  latexmk -pdf CV-Europass-IT.tex
% Versione pubblica: niente telefono, data di nascita o nazionalità. Se un bando
% li richiede, aggiungere \ecvmobile, \ecvdateofbirth e \ecvnationality in una
% copia locale non tracciata (vedi README).
%-------------------------------------------------------------------------------
\documentclass[italian,a4paper]{europasscv}

\usepackage[italian]{babel}
\usepackage[utf8]{inputenc}
\usepackage[T1]{fontenc}
\urlstyle{same}

\ecvname{Mattia Giovannini}
\ecvaddress{Bologna, Italia}
\ecvemail{mattia.giovannini6@studio.unibo.it}
\ecvhomepage{mattia-giovannini.github.io github.com/frigori}

\begin{document}
\begin{europasscv}

\ecvpersonalinfo

\ecvbigitem{Profilo}{Studente della Laurea Magistrale in Ingegneria Informatica (Università di Bologna) in doppio titolo con il Master in Ubiquitous Systems and Networks (Université Côte d'Azur)}

%-------------------------------------------------------------------------------
\ecvsection{Istruzione e formazione}

\ecvtitlelevel{2026 -- 2028 (in corso)}{Laurea Magistrale in Ingegneria Informatica — Doppio Titolo}{}
\ecvitem{}{Università di Bologna, Bologna (Italia) \& Université Côte d'Azur, Polytech Nice Sophia, Nizza (Francia)}
\ecvitem{}{Percorso biennale a doppio titolo con il Master in \textbf{Ubiquitous Systems and Networks}.}

\ecvtitlelevel{2021 -- luglio 2026}{Laurea Triennale in Ingegneria Informatica — voto 94/110}{}
\ecvitem{}{Università di Bologna, Bologna (Italia)}
\ecvitem{}{
  \begin{ecvitemize}
    \item \textbf{Tesi:} \textit{Rilevamento del Disturbo dello Spettro Autistico mediante Machine Learning applicato a Dati di Eye Tracking}
    \item Corsi rilevanti: Analisi Matematica, Algebra Lineare e Geometria, Probabilità e Statistica, Fondamenti di Informatica, Sistemi Operativi, Reti di Calcolatori, Sistemi Informativi, Ingegneria del Software, Tecnologie Web, Architettura degli Elaboratori
  \end{ecvitemize}
}

\ecvtitle{2025}{Progetto di studio Erasmus+}
\ecvitem{}{Université de Tours, Tours (Francia)}
\ecvitem{}{Progetto sul \textbf{machine learning per l'analisi del disturbo dello spettro autistico (ASD)}: classificazione ed estrazione di caratteristiche da dati di eye-tracking, in un contesto accademico internazionale e francofono.}

\ecvtitle{2016 -- 2021}{Diploma di Liceo Scientifico}
\ecvitem{}{Liceo Augusto Righi, Bologna (Italia)}

%-------------------------------------------------------------------------------
\ecvsection{Progetti}

\ecvtitle{2026 (in corso)}{Scheduler Kubernetes resource-aware}
\ecvitem{}{Project work, Enterprise Platforms and DevOps M — Università di Bologna}
\ecvitem{}{Scheduler Kubernetes personalizzato in Python che assegna ogni pod al nodo a latenza minore tra quelli con risorse sufficienti e motiva ogni decisione con eventi Kubernetes; in corso l'estensione allo smart placement di servizi su nodi edge mobili.}

\ecvtitle{2026 (in corso)}{Benchmark di ripristino da guasti di link SDN}
\ecvitem{}{Discrete and Network Optimization Problems M — Università di Bologna \newline \url{https://github.com/frigori/net-restoration-bench}}
\ecvitem{}{Confronto tra backup condiviso (SBPP), Fast Reroute e ricalcolo globale MILP su overhead di capacità, tasso di ripristino e tempo di reazione (Python, NetworkX, PuLP).}

\ecvtitle{2026}{AgenticOS — Audit multi-agente con IA}
\ecvitem{}{Lean Startup Lab, Technology Entrepreneurship \newline \url{https://agenticos-mvp.netlify.app}}
\ecvitem{}{Sistema multi-agente Red Team / Blue Team (Gemini) che analizza perché gli agenti IA di shopping ignorano un prodotto moda e traduce le correzioni in JSON-LD e llms.txt; rilasciato come MVP.}

\ecvtitle{2026}{BOOM GEN AI Challenge 2026 — VEM Sistemi}
\ecvitem{}{Progetto GenAI}
\ecvitem{}{Assistente basato su GenAI per l'analisi automatica di bandi di gara (valutazione GO/NO-GO): modulo di scoring in Python, prompt engineering e Azure OpenAI per l'estrazione dei requisiti.}

\ecvtitle{In corso}{Homelab self-hosted}
\ecvitem{}{Progetto personale \newline \url{https://mattia-giovannini.github.io/homelab/}}
\ecvitem{}{Servizi containerizzati su nodi Raspberry Pi con Docker, Podman rootless e Kubernetes; monitoraggio Prometheus/Grafana e sicurezza CrowdSec.}

%-------------------------------------------------------------------------------
\ecvsection{Esperienze di volontariato e attività}

\ecvtitle{2020 -- 2021}{Webmaster ed Editor Studentesco}
\ecvitem{}{Liceo Augusto Righi, Bologna (Italia)}
\ecvitem{}{Progettazione e gestione delle piattaforme studentesche dell'istituto (\textit{Righi Network} e \textit{Il Sottobanco}).}

\ecvtitle{2016 -- 2020}{Educatore volontario presso centro estivo}
\ecvitem{}{Estate Ragazzi, Bologna (Italia)}
\ecvitem{}{Gestione di gruppi, lavoro di squadra, comunicazione e responsabilità in attività educative.}

%-------------------------------------------------------------------------------
\ecvsection{Competenze personali}

\ecvmothertongue{Italiano}
\ecvlanguageheader
\ecvlanguage{Inglese}{C1}{C1}{C1}{C1}{C1}
\ecvlanguagecertificate{CLA C1, Università di Bologna; Duolingo English Test 130/160 (C1); Cambridge First Certificate B2 (178/190)}
\ecvlastlanguage{Francese}{C1}{C1}{C1}{C1}{C1}
\ecvlanguagecertificate{CLA C1, Università di Bologna; corsi B2/C1 presso l'Université de Tours}
\ecvlanguagefooter

\ecvblueitem{Competenze informatiche}{
  \begin{ecvitemize}
    \item Linguaggi: C, Java, Python, JavaScript, HTML/CSS, SQL, Bash
    \item Machine learning e dati: scikit-learn, NumPy, pandas, Matplotlib, Jupyter Notebook
    \item Sistemi e reti: Linux, TCP/IP, programmazione socket, database relazionali
    \item Strumenti: Git, \LaTeX, MATLAB, Docker, Podman, Kubernetes
  \end{ecvitemize}
}

%-------------------------------------------------------------------------------
\ecvsection{Ulteriori informazioni}

\ecvblueitem{Trattamento dei dati personali}{Autorizzo il trattamento dei miei dati personali ai sensi del D.Lgs. 196/2003 e del Regolamento (UE) 2016/679 (GDPR).}

\end{europasscv}
\end{document}
